% ══════════════════════════════════════════════════════════════════════════════
%  Chapter 16: Desktop Admin Interface (Iced)
% ══════════════════════════════════════════════════════════════════════════════
\chapter{Desktop Admin Interface (Iced)}
\label{ch:desktop-iced}

\begin{learnbox}
\begin{itemize}
  \item Build a desktop admin interface using the Iced framework's Model-View-Update pattern
  \item Implement state management, async API calls, and UI event handling in pure Rust
  \item Design modular UI components with styling and layout
  \item Integrate with the bitcoin-api crate for blockchain communication
\end{itemize}
\end{learnbox}

\lettrine{I}{\index{Iced}ced} is a cross-platform GUI framework for Rust inspired by Elm. It uses a Model-View-Update architecture: your application state lives in a single struct (the \emph{model}), user interactions produce \emph{messages} (a Rust \index{Rust!enum}enum), and a pure \code{update} function handles each message by returning the new state. A separate \code{view} function renders the current state into a widget tree. If you have used Elm, Redux, or SwiftUI, the pattern will be familiar---but \index{Iced}Iced is pure Rust with no JavaScript or web runtime involved.

\begin{notebox}[Prerequisites]
This chapter assumes you have read the Web API chapter (Chapter~\ref{ch:web-api})---the desktop UI is an HTTP client to that \index{REST API}API. No prior GUI framework experience is needed; we introduce \index{Iced}Iced from scratch below. Familiarity with Rust \index{Rust!enum}enums and \index{Rust!pattern matching}pattern matching is essential, as the entire UI is driven by message dispatch.
\end{notebox}

This chapter explains the \code{bitcoin-desktop-ui} crate as a Rust developer reads it: \textbf{how an Iced MVU app boots}, \textbf{how messages flow through the update loop}, and \textbf{how the UI triggers async HTTP calls} to our Rust Bitcoin implementation's admin API.

\section{What This UI Is (In One Sentence)}
\label{sec:ch16-oneline}

The desktop admin UI is a \textbf{thin MVU shell}: it renders state (\code{view.rs}), mutates state in response to messages (\code{update.rs}), and performs \index{Rust!async}async admin calls (\code{api.rs}) on a \index{tokio}Tokio runtime (\code{runtime.rs}).

The methods involved in this architecture are: \code{main} (boot \texttt{+} theme \texttt{+} app start), \code{init\_runtime} and \code{spawn\_on\_tokio} (async boundary), \code{AdminApp::\{new, clear\_related\_data\}} (state and screen cleanup), \code{update} (message dispatcher), and \code{view} and \code{view\_*} helpers (rendering).

% ──────────────────────────────────────────────────────────────────────────────
\section{Diagram: MVU Event Loop (Iced)}
\label{sec:ch16-mvu-loop}

\begin{minted}[fontsize=\footnotesize]{text}
User input (click / type)
   |
   v
Message (enum variant)  -----+
   |                         |
   v                         |
update(&mut AdminApp, Message)
   |
   |-- pure state change -> Task::none()
   |-- async request -> Task::perform() -+
   |                                      |
   v                                      |
view(&AdminApp) -> Element<Message>       |
                                   (later) |
                                Message::XxxLoaded(Result)
\end{minted}

% ──────────────────────────────────────────────────────────────────────────────
\section{Diagram: Async Boundary (HTTP Requests)}
\label{sec:ch16-async-boundary}

\begin{minted}[fontsize=\footnotesize]{text}
update (sync)        runtime (tokio)           api (async)
-------------        ---------------           -----------
Task::perform(   spawn_on_tokio(fut)  ---->  AdminClient::new()
  spawn_on_tokio,     .await                .endpoint().await
  Message::Loaded                          .send()
)
\end{minted}

% ──────────────────────────────────────────────────────────────────────────────
\section{File Organization}
\label{sec:ch16-files}

Together, these chapters walk through implementation in:

\begin{itemize}
  \item \code{bitcoin-desktop-ui-iced/src/main.rs}
  \item \code{bitcoin-desktop-ui-iced/src/runtime.rs}
  \item \code{bitcoin-desktop-ui-iced/src/types.rs}
  \item \code{bitcoin-desktop-ui-iced/src/app.rs}
  \item \code{bitcoin-desktop-ui-iced/src/api.rs}
  \item \code{bitcoin-desktop-ui-iced/src/update.rs}
  \item \code{bitcoin-desktop-ui-iced/src/view.rs}
\end{itemize}

% ──────────────────────────────────────────────────────────────────────────────
\section{Application Entrypoint}
\label{sec:ch16-main}

\begin{listing}[H]
\caption{App entrypoint (\code{bitcoin-desktop-ui-iced/src/main.rs})}
\label{lst:ch16-main}
\begin{minted}[fontsize=\footnotesize]{rust}
mod api;
mod app;
mod runtime;
mod types;
mod update;
mod view;

use app::AdminApp;
use iced::{Theme, application};
use runtime::init_runtime;
use update::update;
use view::view;

fn main() -> iced::Result {
    // Initialize tracing subscriber for logging
    tracing_subscriber::fmt()
        .with_env_filter(tracing_subscriber::EnvFilter::from_default_env())
        .with_target(false)
        .init();

    // Initialize Tokio runtime for async operations
    init_runtime();

    // Run the application
    application(AdminApp::new, update, view)
        .title(title)
        .theme(theme)
        .run()
}
\end{minted}
\end{listing}

\textbf{Walkthrough: what \code{main} wires together:}

\begin{itemize}
  \item \textbf{Modules}: \code{update} and \code{view} are the MVU ``engine''; everything else exists to support them.
  \item \textbf{\code{init\_runtime()}}: starts an async executor so later HTTP tasks don't block the UI thread.
  \item \textbf{\code{application(AdminApp::new, update, view).run()}}: creates initial state once, then enters the event loop:
  \begin{itemize}
    \item \code{view(\&AdminApp)} draws the UI
    \item user interaction emits \code{Message}
    \item \code{update(\&mut AdminApp, Message)} mutates state and optionally schedules work
  \end{itemize}
\end{itemize}

% ──────────────────────────────────────────────────────────────────────────────
\section{Tokio Runtime Bridge}
\label{sec:ch16-tokio-runtime}

\begin{listing}[H]
\caption{Tokio runtime bridge (\code{bitcoin-desktop-ui-iced/src/runtime.rs})}
\label{lst:ch16-runtime}
\begin{minted}[fontsize=\footnotesize]{rust}
use std::sync::OnceLock;

// Store the Tokio runtime handle globally so tasks can access it
static TOKIO_HANDLE: OnceLock<tokio::runtime::Handle> = OnceLock::new();

pub fn init_runtime() {
    // Create a Tokio runtime for async operations
    // This must outlive the application to keep the reactor running
    let rt = tokio::runtime::Runtime::new()
        .expect("Failed to create Tokio runtime");

    // Store the handle globally so it can be accessed from any thread
    TOKIO_HANDLE
        .set(rt.handle().clone())
        .expect("Failed to set Tokio handle");

    // Keep the runtime alive in a background thread
    std::thread::spawn(move || {
        rt.block_on(async {
            // Keep the runtime alive indefinitely
            std::future::pending::<()>().await;
        });
    });
}

// Helper function to wrap a future to ensure it runs on Tokio runtime
pub fn spawn_on_tokio<F>(fut: F)
    -> impl std::future::Future<Output = F::Output> + Send
where
    F: std::future::Future + Send + 'static,
    F::Output: Send + 'static,
{
    let handle = TOKIO_HANDLE
        .get()
        .expect("Tokio runtime not initialized")
        .clone();
    async move { handle.spawn(fut).await.unwrap() }
}
\end{minted}
\end{listing}

\textbf{Walkthrough: why we need \code{spawn\_on\_tokio}:}

Iced's update loop is \textbf{synchronous}: it should be fast and never block on network I/O. The \code{spawn\_on\_tokio(...)} helper gives us a clean ``async boundary'':

\begin{itemize}
  \item \code{update.rs} can return a \code{Task<Message>} immediately (UI stays responsive)
  \item Tokio runs the future in the background
  \item the result is converted into a \code{Message::XxxLoaded(...)} and re-enters the update loop
\end{itemize}

% ──────────────────────────────────────────────────────────────────────────────
\section{Navigation and Events}
\label{sec:ch16-types}

The \code{types.rs} file defines navigation enums and the \code{Message} enum, which is the ``event vocabulary'' of the UI:

\begin{listing}[htbp]
\caption{Navigation enums (\code{bitcoin-desktop-ui-iced/src/types.rs})}
\label{lst:ch16-menu-enum}
\begin{minted}[fontsize=\footnotesize]{rust}
#[derive(Debug, Clone, Copy, PartialEq, Eq)]
pub enum Menu {
    Blockchain,
    Wallet,
    Transactions,
    Mining,
    Health,
}

#[derive(Debug, Clone, Copy, PartialEq, Eq)]
pub enum WalletSection {
    GetWalletInfo,
    GetBalance,
    Create,
    Send,
    History,
    Addresses,
}

#[derive(Debug, Clone, Copy, PartialEq, Eq)]
pub enum BlockchainSection {
    Info,
    LatestBlocks,
    AllBlocks,
    BlockByHash,
}
\end{minted}
\end{listing}

The \code{Message} enum encodes all possible user interactions and async outcomes:

\begin{listing}[htbp]
\caption{Message enum (excerpt) (\code{bitcoin-desktop-ui-iced/src/types.rs})}
\label{lst:ch16-message-enum}
\begin{minted}[fontsize=\footnotesize]{rust}
#[derive(Debug, Clone)]
pub enum Message {
    // Menu/UI navigation
    MenuChanged(Menu),
    WalletSectionChanged(WalletSection),
    BlockchainSectionChanged(BlockchainSection),

    // Configuration inputs
    BaseUrlChanged(String),
    ApiKeyChanged(String),

    // Field inputs
    BlockHashChanged(String),
    MiningAddressChanged(String),
    MiningNBlocksChanged(String),
    TxidChanged(String),

    // Request messages (trigger API calls)
    FetchInfo,
    FetchBlocks,
    FetchBlockByHash(String),
    GenerateToAddress {
        address: String,
        nblocks: u32,
        maxtries: Option<u32>,
    },

    // Result messages (from async API calls)
    InfoLoaded(Result<ApiResponse<BlockchainInfo>, String>),
    BlocksLoaded(Result<ApiResponse<Vec<BlockSummary>>, String>),
    BlockByHashLoaded(Result<ApiResponse<Value>, String>),
    GenerateToAddressDone(Result<ApiResponse<Value>, String>),
}
\end{minted}
\end{listing}

\textbf{Walkthrough: \code{Message} is the UI's event vocabulary:}

Treat \code{Message} as a \textbf{strict interface} between the view and the update loop:

\begin{itemize}
  \item \textbf{Input messages} (\code{BaseUrlChanged}, \code{MiningAddressChanged}, \ldots): pure state changes.
  \item \textbf{Request messages} (\code{FetchInfo}, \code{FetchBlocks}, \code{CreateWalletAdmin}, \ldots): these trigger \code{Task::perform(...)} in the update loop.
  \item \textbf{Result messages} (\code{InfoLoaded}, \code{BlocksLoaded}, \code{TxSent}, \ldots): these store JSON into \code{AdminApp} and update status/editor buffers.
  \item \textbf{UI mechanics} (hover \texttt{+} section routing): \code{*MenuHovered(bool)} and \code{*SectionChanged(...)} exist to drive popups and the active panel.
\end{itemize}

This design keeps the view layer simple: it never ``does'' network calls---it only emits intent.

% ──────────────────────────────────────────────────────────────────────────────
\section{Application State}
\label{sec:ch16-app-state}

\begin{listing}[H]
\caption{Application state (\code{bitcoin-desktop-ui-iced/src/app.rs})}
\label{lst:ch16-adminapp}
\begin{minted}[fontsize=\footnotesize]{rust}
#[derive(Debug)]
pub struct AdminApp {
    pub menu: Menu,
    pub base_url: String,
    pub api_key: String,
    pub status: String,
    pub info: Option<BlockchainInfo>,
    pub blocks: Vec<BlockSummary>,

    // Inputs for actions
    pub block_hash_input: String,
    pub mining_address_input: String,
    pub mining_nblocks_input: String,
    pub txid_input: String,

    // Wallet admin state
    pub wallet_label_input: String,
    pub addresses: Vec<String>,
    pub wallet_info: Option<Value>,
    pub wallet_balance: Option<Value>,
    pub created_wallet_address: Option<String>,

    // Send transaction state
    pub send_from_address: String,
    pub send_to_address: String,
    pub send_amount: String,

    // Navigation state
    pub wallet_section: WalletSection,
    pub transaction_section: TransactionSection,
    pub blockchain_section: BlockchainSection,
    pub mining_section: MiningSection,
    pub health_section: HealthSection,

    // Menu hover state for popups
    pub blockchain_menu_hovered: bool,
    pub wallet_menu_hovered: bool,
    pub transaction_menu_hovered: bool,
    pub mining_menu_hovered: bool,
    pub health_menu_hovered: bool,

    // Response data storage
    pub blocks_all_data: Option<Value>,
    pub block_by_hash_data: Option<Value>,
    pub mining_info_data: Option<Value>,
    pub health_data: Option<Value>,
    pub mempool_data: Option<Value>,
    pub transactions_data: Option<Value>,

    // Text editor states for selectable JSON display
    pub transactions_editor: iced::widget::text_editor::Content,
    pub mempool_editor: iced::widget::text_editor::Content,
    pub wallet_info_editor: iced::widget::text_editor::Content,
    pub blockchain_info_editor: iced::widget::text_editor::Content,
}

impl AdminApp {
    pub fn new() -> (Self, iced::Task<Message>) {
        (
            Self {
                menu: Menu::Blockchain,
                base_url: "http://127.0.0.1:8080".into(),
                api_key: std::env::var("BITCOIN_API_ADMIN_KEY")
                    .unwrap_or_else(|_| "admin-secret".into()),
                status: String::new(),
                // ... (remaining fields initialized)
            },
            iced::Task::none(),
        )
    }

    /// Clear previously loaded data for a specific section
    pub fn clear_related_data(&mut self, section: DataSection) {
        // "Hygiene rule": when we load one dataset, clear adjacent datasets
        // so the UI doesn't render stale JSON next to fresh JSON.
        match section {
            DataSection::BlockchainInfo => {
                self.blocks_all_data = None;
                self.block_by_hash_data = None;
            }
            DataSection::BlocksAll => {
                self.block_by_hash_data = None;
            }
            // ... (more cases)
        }
    }
}
\end{minted}
\end{listing}

\textbf{Walkthrough: what lives in \code{AdminApp} (and why):}

\code{AdminApp} is intentionally ``wide'': it contains both \textbf{domain-ish data} and \textbf{UI presentation state}.

\begin{itemize}
  \item \textbf{Configuration}: \code{base\_url}, \code{api\_key} (what server we talk to).
  \item \textbf{Navigation}: \code{menu} \texttt{+} section enums (which panel to render).
  \item \textbf{Inputs}: text fields like \code{block\_hash\_input}, \code{send\_amount} (what the user typed).
  \item \textbf{Responses}: \code{*\_data: Option<Value>} are cached JSON payloads to display.
  \item \textbf{Editors}: \code{*\_editor: text\_editor::Content} exist so JSON can be selectable/copyable and scrollable.
\end{itemize}

The \code{clear\_related\_data(...)} helper is a small UX rule: when you fetch one dataset, clear nearby datasets so stale JSON doesn't linger on screen.

% ──────────────────────────────────────────────────────────────────────────────
\section{Async API Client Helpers}
\label{sec:ch16-api}

\begin{listing}[H]
\caption{Async API client helpers (\code{bitcoin-desktop-ui-iced/src/api.rs})}
\label{lst:ch16-api}
\begin{minted}[fontsize=\footnotesize]{rust}
use bitcoin_api::{AdminClient, ApiConfig, ApiResponse, BlockchainInfo};

pub async fn fetch_info(cfg: ApiConfig)
    -> Result<ApiResponse<BlockchainInfo>, String>
{
    // A new client is constructed per call; config contains base_url + API key.
    let client = AdminClient::new(cfg).map_err(|e| e.to_string())?;
    client
        .get_blockchain_info()
        .await
        .map_err(|e| e.to_string())
}

pub async fn fetch_blocks(cfg: ApiConfig)
    -> Result<ApiResponse<Vec<BlockSummary>>, String>
{
    let client = AdminClient::new(cfg).map_err(|e| e.to_string())?;
    client.get_latest_blocks().await.map_err(|e| e.to_string())
}

pub async fn fetch_block_by_hash(
    cfg: ApiConfig,
    hash: String,
) -> Result<ApiResponse<Value>, String> {
    let client = AdminClient::new(cfg).map_err(|e| e.to_string())?;
    client
        .get_block_by_hash(&hash)
        .await
        .map_err(|e| e.to_string())
}

pub async fn fetch_mining_info(cfg: ApiConfig)
    -> Result<ApiResponse<Value>, String>
{
    let client = AdminClient::new(cfg).map_err(|e| e.to_string())?;
    client.get_mining_info().await.map_err(|e| e.to_string())
}

pub async fn generate_to_address(
    cfg: ApiConfig,
    address: String,
    nblocks: u32,
    maxtries: Option<u32>,
) -> Result<ApiResponse<Value>, String> {
    let client = AdminClient::new(cfg).map_err(|e| e.to_string())?;
    client
        .generate_to_address(&address, nblocks, maxtries)
        .await
        .map_err(|e| e.to_string())
}

pub async fn fetch_health(cfg: ApiConfig)
    -> Result<ApiResponse<Value>, String>
{
    let client = AdminClient::new(cfg).map_err(|e| e.to_string())?;
    client.health().await.map_err(|e| e.to_string())
}

pub async fn fetch_mempool(cfg: ApiConfig)
    -> Result<ApiResponse<Value>, String>
{
    let client = AdminClient::new(cfg).map_err(|e| e.to_string())?;
    client.get_mempool().await.map_err(|e| e.to_string())
}

pub async fn create_wallet_admin(
    cfg: ApiConfig,
    req: CreateWalletRequest,
) -> Result<ApiResponse<CreateWalletResponse>, String> {
    let client = AdminClient::new(cfg).map_err(|e| e.to_string())?;
    client
        .create_wallet_admin(&req)
        .await
        .map_err(|e| e.to_string())
}

pub async fn send_transaction(
    cfg: ApiConfig,
    req: SendTransactionRequest,
) -> Result<ApiResponse<SendTransactionResponse>, String> {
    let client = AdminClient::new(cfg).map_err(|e| e.to_string())?;
    client
        .send_transaction_admin(&req)
        .await
        .map_err(|e| e.to_string())
}
\end{minted}
\end{listing}

\textbf{Walkthrough: how \code{api.rs} supports the update loop:}

This module is intentionally boring: each function is a \textbf{single HTTP call} wrapped in an \code{async fn} with a UI-friendly error type.

\begin{itemize}
  \item \textbf{Inputs}: \code{ApiConfig} is built from \code{AdminApp \{ base\_url, api\_key \}} in the update function.
  \item \textbf{Return type}: \code{Result<ApiResponse<T>, String>}
  \begin{itemize}
    \item \code{Err(String)} means: ``we could not even complete the request'' (client construction, network error, decode error, etc.).
    \item \code{Ok(ApiResponse<T>)} means: ``we got a response envelope''; it can still represent application-level failure via \code{api.success == false}.
  \end{itemize}
  \item \textbf{Why build a client per call?} It keeps the UI side stateless and makes each request self-contained. In the update loop you can read each \code{Task::perform(spawn\_on\_tokio(api\_call), Message::XxxLoaded)} arm as a direct mapping: \emph{message \(\to\) API call \(\to\) loaded message}.
\end{itemize}

% ──────────────────────────────────────────────────────────────────────────────
\section{The Update Loop: Message Dispatcher}
\label{sec:ch16-update}

\begin{listing}[H]
\caption{Update loop (excerpt) (\code{bitcoin-desktop-ui-iced/src/update.rs})}
\label{lst:ch16-update}
\begin{minted}[fontsize=\footnotesize]{rust}
use crate::api::*;
use crate::app::AdminApp;
use crate::runtime::spawn_on_tokio;
use crate::types::{DataSection, Menu, Message};
use bitcoin_api::ApiConfig;
use iced::Task;

pub fn update(app: &mut AdminApp, message: Message) -> Task<Message> {
    // This match is the "traffic controller" of MVU:
    // - input events update state immediately
    // - button presses spawn async work
    // - async results re-enter as Message::XxxLoaded(...)
    match message {
        // Pure state changes (no I/O)
        Message::MenuChanged(m) => {
            app.menu = m;
            Task::none()
        }
        Message::BaseUrlChanged(v) => {
            app.base_url = v;
            Task::none()
        }
        Message::ApiKeyChanged(v) => {
            app.api_key = v;
            Task::none()
        }
        Message::BlockHashChanged(v) => {
            app.block_hash_input = v;
            Task::none()
        }

        // Async request: spawn API call on Tokio runtime
        Message::FetchInfo => {
            let cfg = ApiConfig {
                base_url: app.base_url.clone(),
                api_key: Some(app.api_key.clone()),
            };
            Task::perform(
                spawn_on_tokio(fetch_info(cfg)),
                Message::InfoLoaded,
            )
        }

        // Result handler: store JSON in app, update UI feedback
        Message::InfoLoaded(res) => {
            match res {
                Ok(api) => {
                    if api.success {
                        app.info = api.data.clone();
                        if let Some(ref data) = api.data {
                            let json_str = serde_json::to_string_pretty(data)
                                .unwrap_or_else(|_| "Error formatting".to_string());
                            app.blockchain_info_editor =
                                iced::widget::text_editor::Content::with_text(&json_str);
                        }
                        app.status = "Blockchain info loaded".into();
                    } else {
                        app.status = api.error
                            .unwrap_or_else(|| "Error loading info".into());
                        app.info = None;
                    }
                }
                Err(e) => {
                    app.status = e;
                    app.info = None;
                }
            }
            Task::none()
        }

        Message::FetchBlocks => {
            app.clear_related_data(DataSection::Blocks);
            let cfg = ApiConfig {
                base_url: app.base_url.clone(),
                api_key: Some(app.api_key.clone()),
            };
            Task::perform(
                spawn_on_tokio(fetch_blocks(cfg)),
                Message::BlocksLoaded,
            )
        }

        Message::BlocksLoaded(res) => {
            match res {
                Ok(api) => {
                    if api.success {
                        app.blocks = api.data.unwrap_or_default();
                        app.status = format!("Loaded {} blocks", app.blocks.len());
                    } else {
                        app.status = api.error
                            .unwrap_or_else(|| "Error loading blocks".into());
                    }
                }
                Err(e) => {
                    app.status = e;
                }
            }
            Task::none()
        }

        // Similar patterns for all other messages...
        _ => Task::none(),
    }
}
\end{minted}
\end{listing}

\textbf{Walkthrough: the async ``round-trip'' pattern (UI \(\to\) HTTP \(\to\) UI):}

The update loop follows a consistent three-step pattern for every API call:

\begin{enumerate}
  \item \textbf{Input message arrives} (user clicks button or types).
  \item \textbf{Update handler} (sync):
  \begin{itemize}
    \item Mutates local state immediately (inputs/status/hover state).
    \item Returns \code{Task::perform(spawn\_on\_tokio(async \{ http\_call().await \}), Message::XxxLoaded)}.
  \end{itemize}
  \item \textbf{Async work happens} (on Tokio runtime, background):
  \begin{itemize}
    \item \code{api.rs} \(\to\) \code{AdminClient} \(\to\) HTTP \(\to\) JSON.
  \end{itemize}
  \item \textbf{Result message arrives} (\code{Message::XxxLoaded(Result<ApiResponse<T>, String>)}).
  \item \textbf{Update handler} (sync, again):
  \begin{itemize}
    \item Stores JSON payload in \code{app.*\_data}.
    \item Updates text editor buffers (pretty JSON).
    \item Updates \code{app.status} for UX feedback.
  \end{itemize}
\end{enumerate}

This separates concerns cleanly: the view never blocks, \code{update} remains lightweight, and I/O happens on the Tokio runtime.

% ──────────────────────────────────────────────────────────────────────────────
\section{The View Layer: Rendering and Dispatch}
\label{sec:ch16-view}

\begin{listing}[H]
\caption{View function (excerpt) (\code{bitcoin-desktop-ui-iced/src/view.rs})}
\label{lst:ch16-view}
\begin{minted}[fontsize=\footnotesize]{rust}
use crate::app::AdminApp;
use crate::types::{Menu, Message, BlockchainSection};
use iced::Element;
use iced::widget::{button, column, container, row, text, text_input};

pub fn view<'a>(app: &'a AdminApp) -> Element<'a, Message> {
    // The view is (mostly) pure: it reads app and returns an Element<Message>.
    // Any "actions" are triggered by Messages, handled in update.rs.

    // Toolbar: configure base_url and api_key
    let toolbar = row![
        text("Base URL:").size(12),
        text_input("http://127.0.0.1:8080", &app.base_url)
            .on_input(Message::BaseUrlChanged)
            .width(iced::Length::Fixed(250.0)),
        text("Admin Key:").size(12),
        text_input("admin-secret", &app.api_key)
            .on_input(Message::ApiKeyChanged)
            .width(iced::Length::Fixed(200.0)),
        text(format!("Status: {}", app.status))
            .size(12)
            .width(iced::Length::Fill),
    ]
    .spacing(8)
    .padding(8);

    // Left navigation menu
    let nav_buttons = column![
        button("Blockchain")
            .on_press(Message::MenuChanged(Menu::Blockchain)),
        button("Wallet")
            .on_press(Message::MenuChanged(Menu::Wallet)),
        button("Transactions")
            .on_press(Message::MenuChanged(Menu::Transactions)),
        button("Mining")
            .on_press(Message::MenuChanged(Menu::Mining)),
        button("Health")
            .on_press(Message::MenuChanged(Menu::Health)),
    ]
    .spacing(4)
    .padding(8);

    // Main content panel: dispatch to view_* functions by menu
    let content = match app.menu {
        Menu::Blockchain => view_blockchain(app),
        Menu::Wallet => view_wallet(app),
        Menu::Transactions => view_transactions(app),
        Menu::Mining => view_mining(app),
        Menu::Health => view_health(app),
    };

    // Assemble: toolbar + (nav + content)
    container(
        column![
            toolbar,
            row![
                container(nav_buttons).width(iced::Length::Fixed(150.0)),
                container(content).width(iced::Length::Fill),
            ]
            .spacing(8)
            .height(iced::Length::Fill),
        ]
        .spacing(8)
    )
    .padding(8)
    .into()
}

// Helper: render the Blockchain section
fn view_blockchain<'a>(app: &'a AdminApp) -> Element<'a, Message> {
    match app.blockchain_section {
        BlockchainSection::Info => {
            column![
                button("Fetch Info").on_press(Message::FetchInfo),
                text_editor(&app.blockchain_info_editor)
                    .on_action(Message::BlockchainInfoEditorAction),
            ]
            .spacing(8)
            .into()
        }
        BlockchainSection::LatestBlocks => {
            column![
                button("Fetch Blocks").on_press(Message::FetchBlocks),
                row![
                    text(format!("Loaded: {} blocks", app.blocks.len())),
                ]
                .spacing(8),
            ]
            .spacing(8)
            .into()
        }
        _ => container(text("Select a section")).into(),
    }
}

// Similar for view_wallet, view_transactions, view_mining, view_health...
\end{minted}
\end{listing}

\textbf{Walkthrough: view dispatch \texttt{+} popup hover model:}

The view function is organized into layers:

\begin{enumerate}
  \item \textbf{Toolbar} (top): configuration inputs (\code{base\_url}, \code{api\_key}) and status feedback.
  \item \textbf{Navigation} (left): menu buttons dispatch \code{Message::MenuChanged(...)}.
  \item \textbf{Main content panel} (right): matches on \code{app.menu} to dispatch to \code{view\_blockchain}, \code{view\_wallet}, etc.
\end{enumerate}

Each content panel handles its own sub-sections and data display:

\begin{itemize}
  \item Buttons emit request messages (\code{FetchInfo}, \code{FetchBlocks}, etc.).
  \item Text editors display and allow copying of JSON responses.
  \item Input fields emit change messages (\code{BlockHashChanged}, etc.).
  \item Hover states drive popup menus: when \code{app.blockchain\_menu\_hovered == true}, the UI renders a submenu; when false, it disappears.
\end{itemize}

The view layer never calls the API directly---all side effects go through messages and the update loop.

% ──────────────────────────────────────────────────────────────────────────────
\begin{chaptersummary}
\item We built a full-featured desktop admin interface using Iced's Model-View-Update pattern in pure Rust.
\item We implemented the complete MVU cycle: Model state (\code{AdminApp}), Message types, Update logic (\code{update} function), and View rendering (\code{view} function).
\item We integrated async API communication with the bitcoin-api crate for real-time blockchain data via a Tokio runtime bridge.
\item We designed modular UI components with popup menus, selectable text, and a consistent styling system.
\item The async boundary between sync MVU and Tokio-powered HTTP keeps the UI responsive while allowing long-running network calls.
\end{chaptersummary}

% ──────────────────────────────────────────────────────────────────────────────
\section{Exercises}
\label{sec:ch16-exercises}

\begin{enumerate}
  \item \textbf{Add a Refresh Button} --- Add a ``Refresh'' button to the admin UI that reloads blockchain data from the API. Define a new \code{Message::Refresh} variant, implement the update handler to trigger an API call, and add the button to the view. Observe how the MVU cycle processes this new interaction.

  \item \textbf{MVU Cycle Trace} --- Pick any user interaction (e.g., clicking ``Mine Block'') and trace the complete MVU cycle: which \code{Message} variant is generated, what the \code{Update} function does, what side effects (API calls, state changes) occur, and how the \code{View} reflects the new state.
\end{enumerate}

% ──────────────────────────────────────────────────────────────────────────────
\section{Further Reading}
\label{sec:ch16-reading}

\begin{itemize}
  \item \textbf{Iced Framework Documentation} --- API reference for the Iced GUI framework at \url{https://docs.rs/iced/}.
  \item \textbf{Iced Examples Repository} --- Official examples demonstrating various Iced patterns at \url{https://github.com/iced-rs/iced/tree/master/examples}.
  \item \textbf{The Elm Architecture} --- The original MVU pattern that inspired Iced's design at \url{https://guide.elm-lang.org/architecture/}.
\end{itemize}

